\subsection{Inteligencia artificial en la educación}

La integración de la inteligencia artificial (IA) en la educación ha sido objeto de creciente investigación en los últimos años. Esta tecnología emergente ofrece oportunidades significativas para transformar las prácticas pedagógicas y mejorar los resultados de aprendizaje de los estudiantes.

\cite{guerreirosantallaDesarrolloPlanEstudios2023} desarrollaron un plan de estudios para la enseñanza de inteligencia artificial en educación preuniversitaria, destacando la importancia de familiarizar a los estudiantes con conceptos fundamentales de IA desde etapas tempranas. Este trabajo evidencia la necesidad de preparar a las nuevas generaciones para un mundo donde la IA será omnipresente en diversos ámbitos profesionales y cotidianos.

Por su parte, \cite{panquebanInteligenciaArtificialEducacion2024a} realizaron una revisión sistemática sobre la aplicación de la inteligencia artificial en la enseñanza de las matemáticas, identificando tendencias actuales y desafíos en la implementación de estas herramientas tecnológicas. Los autores señalan que la IA puede personalizar la experiencia de aprendizaje, adaptándose al ritmo y estilo de cada estudiante, lo que representa un avance significativo respecto a los métodos tradicionales de enseñanza.

Estos estudios muestran que la IA en educación no solo es una herramienta tecnológica, sino también un medio para repensar las prácticas pedagógicas y promover un aprendizaje más inclusivo y efectivo.
